\documentclass[journal]{IEEEtran}
\usepackage{graphicx}
\usepackage{amsmath,amssymb}
\usepackage{booktabs}
\usepackage{xcolor}
\usepackage[hidelinks]{hyperref}
\usepackage{url}
\Urlmuskip=0mu plus 1mu\relax
\makeatletter
\g@addto@macro\UrlBreaks{\do\a\do\b\do\c\do\d\do\e\do\f\do\g\do\h\do\i\do\j\do\k\do\l\do\m\do\n\do\o\do\p\do\q\do\r\do\s\do\t\do\u\do\v\do\w\do\x\do\y\do\z\do\A\do\B\do\C\do\D\do\E\do\F\do\G\do\H\do\I\do\J\do\K\do\L\do\M\do\N\do\O\do\P\do\Q\do\R\do\S\do\T\do\U\do\V\do\W\do\X\do\Y\do\Z\do\0\do\1\do\2\do\3\do\4\do\5\do\6\do\7\do\8\do\9\do\-\do\.\do\_\do\?\do\=\do\&}
\makeatother

\begin{document}

\title{A Decision-Focused, Fairness-Constrained Framework for Multi-Horizon Graft-Rejection Prediction and Equitable Kidney Allocation in End-Stage Renal Disease}

\author{Author~One,~Author~Two,~and~Corresponding~Author%
\thanks{Manuscript received November 2026. The authors are with the Department/Institute, City, Country (e-mail: corresponding@institute.edu).}%
\thanks{\emph{Code and Data Availability:} all code, the synthetic generator, and the experiment scripts are at \url{https://github.com/sarthaksaptarshi/kidney-allocation-dfl}, with pinned dependencies (\texttt{requirements.txt}) and a reproduction map (\texttt{REPRODUCE.md}) linking each figure and table to the script that generates it. The open UCI, Mendeley (CC BY 4.0), and Jakarta (Dryad) datasets are obtained from their cited sources.}}

\markboth{IEEE Transactions on Artificial Intelligence, Draft, 2026}%
{Author \MakeLowercase{\textit{et al.}}: Decision-Focused, Fairness-Constrained Kidney Graft Prediction and Allocation}

\maketitle

\begin{abstract}
Kidney transplantation is the preferred therapy for end-stage renal disease, yet rejection prediction and allocation of scarce kidneys are solved separately: a risk model trained for accuracy feeds a downstream allocator. We propose a decision-focused, fairness-constrained framework that predicts graft-rejection risk across horizons with a censoring-aware, temporally consistent survival model; attaches conformal uncertainty for biopsy abstention, and separates antibody-mediated rejection from tacrolimus toxicity; and feeds a surrogate benefit into a fairness-constrained optimizer with a NOTTO-inspired equity priority, closing a decision-focused loop. On synthetic data ($n=18{,}000$) the offer-time predictor matches Cox at every horizon (AUROC $0.825$ vs.\ $0.822$) while temporally consistent; uncertainty gating lifts rejection-versus-toxicity accuracy to $0.86$. The central finding is negative: decision-focused training lowers regret only for a misspecified predictor in the unconstrained assignment; correcting the features, or activating equity or ABO/crossmatch feasibility, leaves it indistinguishable from zero---a two-stage predictor matches it, so feasibility, equity, and calibration govern allocation quality. The equity priority raises \emph{simulated} access for women, sensitized, and pediatric candidates at a $0.85\%$ fairness price. Across three open datasets Cox leads at $31$ events but ties on a larger cohort; allocation and fairness results remain simulated.
\end{abstract}

\begin{IEEEkeywords}
Decision-focused learning, survival analysis, organ allocation, algorithmic fairness, conformal prediction, kidney transplantation.
\end{IEEEkeywords}

\IEEEpeerreviewmaketitle

\begin{center}
\fbox{\parbox{0.94\columnwidth}{\footnotesize\textit{Reproducibility note.} The main synthetic results come from one seed-fixed run (seed 20261116, $n=18{,}000$); the decision-focused, integration, fairness, and event-count analyses use repeated runs (ten seeds), stated in situ; all synthetic point estimates are seed-fixed and reproduce across BLAS builds to within $\approx0.03$ regret points, with significance verdicts unchanged. The generator is parameterized from published SRTR/UNOS distributions and NOTTO categories; results show the pipeline is well posed, not that it achieves these numbers clinically. The real-data arms use the openly licensed UCI, Mendeley, and Jakarta datasets.}}
\end{center}
\vspace{4pt}

\section{Introduction}
\IEEEPARstart{K}{idney} transplantation offers longer survival and better quality of life than dialysis for end-stage renal disease, but donor organs are scarce and every avoidable graft loss sends a patient back to the waiting list. Two decisions dominate outcomes: predicting which transplants will reject, and deciding which candidate should receive each available kidney. The machine-learning literature has largely treated these as separate problems, and chained them na\"ively---a risk score trained for accuracy is passed to an allocation rule, even though accuracy and allocation quality are different objectives. This \emph{predict-then-optimize} decoupling is the central gap we address.

\subsection{Literature Survey}
\textbf{Graft-survival and failure prediction.} The largest body of work predicts graft survival from registry data---deep survival networks on the SRTR~\cite{luck}, gradient boosting and random forests on national cohorts~\cite{bae,thong}, and recent multimodal models fusing EHR streams with clinical text (graft-loss AUC $0.96$, rejection AUC $0.84$ on $3{,}382$ recipients)~\cite{tfn}. Systematic reviews temper this: a 2023 meta-analysis found a pooled AUROC near $0.82$~\cite{meta23}, and comparative reviews find no evidence that machine learning beats Cox once both have the same variables~\cite{klundert,bae,dibello}. We take that caution seriously, reporting honest, modest margins rather than headline AUROC.

\textbf{Rejection diagnosis.} A parallel line targets the diagnosis of rejection itself. Reviews catalogue machine-learning approaches to allograft rejection and its personalized management~\cite{belcic,almoussawy}, and deep networks now classify and subtype T-cell- and antibody-mediated rejection directly from whole-slide biopsy images~\cite{wsi}. These operate post-biopsy, downstream of the allocation decision; our differential instead flags the AMR-versus-toxicity ambiguity early and refers uncertain cases for exactly such a biopsy.

\textbf{Tacrolimus management.} Tacrolimus exposure, the dominant modifiable driver of both rejection and nephrotoxicity, has its own literature, from exposure prediction to genetic-algorithm and deep-forest dosing~\cite{tacro}. These optimize drug levels but do not connect them to the allocation decision.

\textbf{Allocation and fairness.} Allocation and exchange form the third strand. Optimization and machine learning have been used to design fairer allocation policies~\cite{bert}, and fairness has been formalized for kidney paired donation through calibration constraints~\cite{zhang} and learned exchange weights~\cite{carvalho}. Others have benchmarked AI predictions of long-term allograft failure against transplant physicians of varying experience~\cite{divard}. Yet these treat prediction and allocation separately, or address paired exchange rather than whole-pool deceased-donor assignment, and none couple a calibrated multi-horizon predictor to the optimizer.

\textbf{The Indian setting.} India is the least-studied setting. Despite the third-highest transplant volume worldwide, it has only one deep-learning allocation prototype on NOTTO guidelines~\cite{reva} and a readiness analysis for unified allocation~\cite{oneindia}. The 2023 ``One Nation, One Organ Allocation'' reform~\cite{notto} and the 2025 equity advisory adding points for women and near relatives of past donors~\cite{nottoadv} create both the need and the opportunity for allocation AI designed for the Indian pool. Table~\ref{tab:related} positions representative studies; no prior strand couples a multi-horizon, uncertainty-aware predictor to a fairness-constrained optimizer.

\begin{table*}[!t]
\centering
\caption{Positioning Against Representative Prior Work. A Check Mark Denotes the Capability Is Present.}
\label{tab:related}
\begin{tabular}{@{}p{3.6cm}p{3.1cm}p{3.4cm}cccc@{}}
\toprule
Study (representative) & Focus & Method & Multi-horizon & Calibrated unc. & Fairness & Joint pred.+alloc. \\
\midrule
Deep graft survival, SRTR~\cite{luck} & Graft-failure survival & Deep survival net & -- & -- & -- & -- \\
Bae \textit{et al.}, 2020~\cite{bae} & DGF, 1-yr AR, failure & GBM / RF vs.\ regression & -- & -- & -- & -- \\
Temporal Fusion Nexus~\cite{tfn} & Graft loss, rejection, death & Multimodal deep net & partial & -- & -- & -- \\
Thongprayoon \textit{et al.}~\cite{thong} & Acute rejection (UNOS) & Random forest & -- & -- & -- & -- \\
Offei \textit{et al.}, 2025~\cite{offei} & Eplet/EMS survival & ML survival & -- & -- & -- & -- \\
Tacrolimus deep forest~\cite{tacro} & Tacrolimus dosing & GA + deep forest & -- & -- & -- & -- \\
Zhang, Dai \& Li, 2025~\cite{zhang} & Kidney paired donation & IP + calibration fairness & -- & -- & \checkmark & partial \\
Carvalho \textit{et al.}~\cite{carvalho} & Kidney exchange weights & RL / learned weights & -- & -- & \checkmark & partial \\
Bertsimas \& Trichakis~\cite{bert} & Allocation policy & ML + optimization & -- & -- & \checkmark & \checkmark (not DF) \\
REVA model (India)~\cite{reva} & NOTTO allocation & Deep neural net & -- & -- & -- & alloc.\ only \\
\textbf{This work} & \textbf{Rejection + allocation} & \textbf{Survival+conformal+fair opt.+DF} & \checkmark & \checkmark & \checkmark & \checkmark (surrogate) \\
\bottomrule
\end{tabular}
\end{table*}

\subsection{Remaining Gaps and Contributions}
Three gaps persist. First, rejection is usually modeled at a single fixed horizon, so the early-AMR risk and the slow chronic injury that governs ten-year survival are never reconciled in one temporally consistent model. Second, most models output a bare probability with no calibrated uncertainty, and cannot separate AMR from the tacrolimus nephrotoxicity that mimics it. Third, allocation equity is treated as an afterthought rather than a constraint, despite documented disparities by sex, sensitization, and socioeconomic status.

We propose one framework that treats prediction and equitable allocation as a single decision problem. Our contributions are: (1) a \emph{multi-horizon, censoring-aware survival predictor} whose cumulative-incidence risk at one-month, one-year, and ten-year horizons is monotone by construction and unbiased under the covariate-dependent censoring that biases na\"ive fixed-horizon classifiers; (2) \emph{calibrated uncertainty and a clinical differential}---conformal abstention that refers ambiguous cases for biopsy, and a gated classifier separating AMR from tacrolimus toxicity; (3) \emph{fairness-constrained allocation with a quantified price of fairness}, encoding NOTTO's reformed bonus structure as soft optimizer priorities; and (4) a \emph{decision-focused loop tested end-to-end}---we train the benefit predictor with an SPO+ objective and first show, as a controlled allocation-module experiment on held-out donor batches with seed-clustered confidence intervals, that it reduces regret versus a matched two-stage predictor in the \emph{utility-only} assignment---but only when the predictor is misspecified, since correcting the feature map already erases the gain; we then close the loop, placing the NOTTO equity priority and the ABO/crossmatch feasible-edge set inside \emph{both} training and evaluation, and a feasibility$\times$priority factorial leaves the integrated paired reduction statistically indistinguishable from zero for this linear learner, while a nonlinear-capacity check and a reduced-form survival-to-allocation study probe---without settling---how far that result extends. The contribution is not SPO+ itself but this transferable finding: once feasibility and equity constraints are active, the realized allocation is robust to how benefit is predicted, so a simple two-stage predictor suffices and effort is better spent on the feasibility and equity layers and on calibration than on the decision-focused objective (Section~\ref{sec:discussion}). All four are demonstrated on a synthetic cohort and corroborated on three openly licensed real-world datasets.

\section{Methodology and Simulation Framework}
This section presents the problem, the four methodological components (Fig.~\ref{fig:arch}), and the simulation framework on which the method is evaluated. The two are described together because each modeling choice is matched by a corresponding property of the data generator, which is what makes the evaluation a fair test of that choice.

\subsection{Problem Formulation}
Let each donor--recipient pairing carry features $x$. The prediction task estimates the cumulative incidence of immune/chronic graft rejection $F(t\mid x)$ at horizons $t\in\{1\text{ month},1\text{ year},10\text{ years}\}$ under right-censoring, plus a differential label (AMR vs.\ tacrolimus toxicity) for early dysfunction. The allocation task, given a batch of $D$ donor kidneys and $R$ waitlisted candidates with a benefit matrix $B$---in our experiments a fixed surrogate of pairwise features, with fitted survival a \emph{proposed} input to a future integrated allocator---selects an assignment maximizing total benefit subject to equity priorities.

\begin{figure*}[!t]
\centering
\includegraphics[width=\textwidth]{out/fig_architecture.png}
\caption{The decision-focused, fairness-constrained pipeline. (A) Multi-horizon survival predictor; (B) conformal uncertainty with an AMR-versus-toxicity gate that abstains to biopsy; (C) fair, feasibility-constrained allocation optimizer (NOTTO priority, ABO/crossmatch). The dashed arrow (D) is the decision-focused, regret-aware feedback that trains the \emph{benefit-matrix learner}---a surrogate in the main experiments, or the reduced-form survival proxy of Section~\ref{sec:loop}---not the survival predictor (A). Conformal abstention (B) is a separate, not-yet-integrated stage.}
\label{fig:arch}
\end{figure*}

\subsection{Multi-Horizon, Censoring-Aware Survival Predictor}
Rather than train one binary classifier per horizon---the common practice, which labels patients censored before the horizon as non-events and is biased under covariate-dependent censoring---we discretize follow-up into intervals and fit a conditional hazard $\hat h_k(x)$ on the patients still \emph{at risk} at the start of interval $k$. Cumulative incidence follows from the product of interval survivals,
\begin{equation}
\hat F(t\mid x)=1-\prod_{k:\,\tau_k\le t}\big(1-\hat h_k(x)\big).
\end{equation}
Because each hazard conditions only on the at-risk set, the estimator is unbiased when censoring is independent of the event given $x$. The construction is monotone across horizons by design; an isotonic projection guarantees it exactly,
\begin{equation}
\hat F(t_1\mid x)\le\hat F(t_2\mid x)\le\hat F(t_3\mid x),\quad t_1<t_2<t_3 .
\end{equation}

\subsection{Calibrated Uncertainty and the AMR/Toxicity Differential}
We recalibrate the one-year risk by isotonic regression and build split-conformal prediction sets at level $\alpha=0.10$. With nonconformity score $s=1-p(\text{true label})$, the threshold is the $(1-\alpha)$ empirical quantile on a calibration split, and the prediction set keeps every admissible label,
\begin{equation}
\hat q=\mathrm{Quantile}_{1-\alpha}\big(\{s_i\}_{i\in\mathrm{cal}}\big),\quad
\mathcal{C}(x)=\{y:1-p_y(x)\le\hat q\}.
\end{equation}
Singleton sets are confident predictions; size-two sets abstain and are referred for biopsy. A separate gradient-boosted classifier, trained only on the early-dysfunction cohort, separates AMR from tacrolimus toxicity, with a confidence-margin gate that refers borderline cases.

\subsection{Fairness-Constrained Allocation}
The optimizer consumes a benefit matrix $B(j,i)$ for donor $j$ and recipient $i$. In the experiments this benefit is a \emph{fixed surrogate}: an interpretable, reproducible function of pairwise features standing in for a future survival-derived benefit, rather than the fitted survival probabilities themselves (end-to-end substitution of the fitted model is future work). We define it in $[0.05,0.98]$,
\begin{equation}
\begin{aligned}
B(j,i)=\mathrm{clip}\big(&\beta_0-\beta_1|h_j-h_i|-\beta_2\,\mathrm{cPRA}_i\\[-1pt]
&-\beta_3(a_j-40)_+-\beta_4(a_i-55)_++\beta_5 e^{-v_i/6}\big),
\end{aligned}
\end{equation}
decreasing in HLA mismatch $|h_j-h_i|$, sensitization $\mathrm{cPRA}_i$, donor age $a_j$ and recipient age $a_i$, and rising with shorter dialysis vintage $v_i$ (with $\beta=(0.90,0.045,0.0018,0.003,0.002,0.01)$ and $(\cdot)_+$ the positive part). Equity enters as a \emph{NOTTO-inspired} additive priority score $s_i$ (weights illustrative of the policy structure, not the gazetted values), combining a dialysis-time credit that rewards \emph{longer} waits $v_i$ with group bonuses, pediatric the largest,
\begin{equation}
s_i \;\propto\; 0.3\,v_i + 2\cdot\mathbb{1}[\text{female}]+5\cdot\mathbb{1}[\text{pediatric}]+4\cdot\mathbb{1}[\text{cPRA}_i>80],
\end{equation}
normalized to $[0,1]$ across the pool. The assignment maximizes benefit plus the weighted priority,
\begin{equation}
\max_{a\in\mathcal A}\ \sum_{j,i}a_{ji}\big(B_{ji}+\lambda\,s_i\big)\ \text{s.t.}\ \textstyle\sum_i a_{ji}\le1,\ \sum_j a_{ji}\le1 .
\end{equation}
These are \emph{soft} priorities, not hard reserved minima: a large pediatric weight approximates NOTTO's automatic pediatric priority without enforcing it as a constraint. Sweeping $\lambda$ traces the equity--utility frontier; the \emph{price of fairness} is the relative utility lost at the chosen operating point versus pure utility maximization. Note the deliberate opposition between $B$ (which rewards \emph{short} vintage, the better medical outcome) and the $s_i$ dialysis credit (which rewards \emph{long} waits): that tension is what the price of fairness measures.

\subsection{Decision-Focused Coupling}
The two-stage pipeline trains the predictor to minimize estimation error, not allocation regret, which we define against the oracle assignment,
\begin{equation}
\mathrm{Regret}(\hat B)=\frac{U^\star-\sum_{(j,i)\in a(\hat B)}B_{ji}}{U^\star},\quad
U^\star=\!\!\sum_{(j,i)\in a(B)}\!\!B_{ji} .
\end{equation}
For the global assignment, regret increased monotonically with benefit-estimation error \emph{in our noise sweep} (Section~\ref{sec:results}-D; we claim this for the observed sweep, not as a general law), so a decision-aware objective---training the predictor to minimize regret directly, which we instantiate with the SPO+ loss~\cite{spo}---can reduce regret even when raw accuracy is unchanged (Section~\ref{sec:results}-D). The loop also supports continuous and federated updating across centers, keeping the model current without moving patient data (relevant to India's DPDP Act, 2023).

To test the loop as a whole rather than the predictor in isolation, we also define the \emph{integrated} objective the full framework actually optimizes. Over the ABO/crossmatch-feasible edge set $\mathcal{F}$, the assignment maximizes benefit plus the normalized NOTTO additive priority $s_i$,
\begin{equation}
a^\star=\arg\max_{a\subseteq\mathcal{F}}\sum_{(j,i)\in a}\big(B_{ji}+\lambda\,s_i\big),
\end{equation}
with $\lambda$ the equity weight, and the \emph{integrated regret} replaces the oracle assignment above with this fair, feasible oracle. Section~\ref{sec:results}-D trains and evaluates the two-stage and SPO+ predictors against this integrated regret---with $s_i$ and $\mathcal{F}$ present in \emph{both} arms---so the comparison speaks to fair, feasible allocation rather than to the utility-only module alone.

\subsection{Public Data for the Real Study}
The framework targets openly obtainable data (Table~\ref{tab:data}). Patient-level Indian transplant data is not openly released, so the allocation pool is simulated from published NOTTO criteria and SRTR distributions---itself a contribution given the data gap.

\begin{table}[!t]
\centering
\caption{Public Data Sources: Used in This Study and Intended for Future Work.}
\label{tab:data}
\begin{tabular}{@{}p{2.1cm}p{1.8cm}p{2.9cm}@{}}
\toprule
Dataset & Access & Role \\
\midrule
UCI-CKD (Apollo) & Fully open & CKD/ESRD detection (used) \\
Mendeley allograft & CC BY 4.0 & Rejection; graft-loss survival (used) \\
Jakarta cohort (Dryad) & Open & Second graft-loss cohort (used) \\
SRTR / OPTN STAR & Free, DUA & Registry re-estimation (future) \\
\bottomrule
\end{tabular}
\end{table}

\subsection{Synthetic Cohort and Generative Model}
We generate $n=18{,}000$ pairings with thirteen features spanning recipient, donor, immunological, and pharmacological axes (Table~\ref{tab:feat}). A shared latent frailty $u\sim\mathcal{N}(0,1)$ encodes unobserved immunological risk and couples the horizons---mirroring the frailty that a single-horizon model cannot see but the multi-horizon estimator exploits. Interval hazards and an informative (covariate-dependent) censoring time are drawn as
\begin{align}
h_1(x)&=\sigma(-2.95+0.9\,\mathrm{imm}(x)+0.8\,u),\\
h_2(x)&=\sigma(-1.70+0.8\,\mathrm{imm}(x)+0.3\,\mathrm{chr}(x)+0.9\,u),\\
h_3(x)&=\sigma(-0.35+0.3\,\mathrm{imm}(x)+0.9\,\mathrm{chr}(x)+1.0\,u),\\
C_i&=\min\!\big(1+\mathrm{Exp}(\lambda_i),120\big),\;\lambda_i=54\,e^{-0.35\,\mathrm{chr}^{\mathrm{std}}(x_i)},
\end{align}
for intervals month 0--1, 1--12, and year 1--10 respectively. The risk scores $\mathrm{imm}(x)$ and $\mathrm{chr}(x)$ are nonlinear, with an HLA$\times$DSA interaction, a sensitization threshold, and a U-shaped tacrolimus-exposure effect, so that tree and survival models have genuine structure to exploit. Critically, sicker (higher-chronic-risk) patients are censored earlier: this covariate-dependent censoring is precisely what biases na\"ive fixed-horizon labels and what the at-risk survival estimator is designed to withstand. The resulting true prevalences are $12.2\%$ (1 month), $28.0\%$ (1 year), and $53.3\%$ (10 years), with $0\%$, $15.2\%$, and $57.9\%$ of patients censored before each horizon (Fig.~\ref{fig:cohort}). A separate dysfunction cohort ($n=4{,}711$) carries the AMR-versus-toxicity label, with toxicity prevalence $36\%$. The allocation pool is a batch of $200$ donor kidneys and $300$ waitlisted candidates, with the surrogate benefit a function of HLA match, cold ischemia, sensitization, and dialysis vintage.

\begin{table}[!t]
\centering
\caption{Feature Groups.}
\label{tab:feat}
\begin{tabular}{@{}p{1.8cm}p{5.6cm}@{}}
\toprule
Group & Variables \\
\midrule
Recipient & age, BMI, dialysis vintage, diabetes, sex \\
Donor & age, deceased-vs-living \\
Immunological & HLA mismatches (0--6), cPRA, DSA surrogate, cold ischemia time \\
Pharmacological & tacrolimus trough, CYP3A5 expresser status \\
\bottomrule
\end{tabular}
\end{table}

\begin{figure}[!t]
\centering
\includegraphics[width=0.74\columnwidth]{out/fig6_cohort.png}
\caption{Synthetic cohort: true rejection prevalence and the fraction censored before each horizon. Censoring grows from $0\%$ at one month to $58\%$ at ten years.}
\label{fig:cohort}
\end{figure}

\subsection{Feature Timing and Leakage Control}
Each predictor is used only at or before the decision that consumes it (Table~\ref{tab:timing}). Recipient demographics, dialysis vintage, and CYP3A5 genotype are fixed at listing; donor factors, HLA mismatch, cPRA, DSA, and crossmatch at the offer. Cold-ischemia time is \emph{estimated} and the induction regimen the \emph{planned} one---both fixed at offer, so both are admissible. Tacrolimus trough is measured only after transplantation, so it enters \emph{only} the AMR-versus-toxicity differential, not the pre-offer models. For the real Mendeley model only pre-transplant variables enter; every post-transplant trajectory variable is withheld. This separation of planned from observed values makes the reported discrimination a forecast, not a retrospective fit.

\begin{table}[!t]
\centering
\caption{Feature Timing by Decision Point. Planned Values Are Distinguished From Observed; Post-Outcome Variables Are Withheld.}
\label{tab:timing}
\begin{tabular}{@{}p{3.0cm}p{2.3cm}p{2.0cm}@{}}
\toprule
Feature group & Known by & Status \\
\midrule
Recipient demographics, dialysis vintage, CYP3A5 genotype & Listing & Used \\
Donor age/type, HLA, cPRA, DSA, crossmatch, cold ischemia & Match/offer & Used \\
Induction regimen (planned) & Offer (planned) & Used \\
Tacrolimus trough & Post-transplant & Differential only \\
Graft function, infections, rejection, failure & Post-transplant & Withheld (leakage) \\
\bottomrule
\end{tabular}
\end{table}

\section{Experimental Results}
\label{sec:results}
Unless stated otherwise, results come from one seed-fixed synthetic run (seed 20261116; 70/30 train--test split; evaluation restricted to patients with known status at each horizon). Three analyses use repeated runs and are labeled where they appear: the decision-focused comparison (ten seeds, $240$ held-out batches), the fairness benchmark (eight independent pools per scarcity regime), and the event-count crossover (nested training subsamples).

\subsection{Multi-Horizon Discrimination and Temporal Consistency}
All horizon models use the \emph{offer-time} feature set (measured tacrolimus excluded; Section~\ref{sec:results}-G). On discrimination the proposed survival model is indistinguishable from the best single-horizon baselines and from Cox at every horizon; at ten years it holds the highest point estimate ($0.825$) but the paired bootstrap against Cox includes zero (Proposed$-$Cox $+0.003$, $95\%$ CI $[-0.009,0.015]$; Fig.~\ref{fig:auroc}a, Table~\ref{tab:auroc}), and re-scoring on the full counterfactual labels gives the same null---no advantage, and no claim of equivalence. The gains lie elsewhere. The model is temporally monotone by construction: \emph{independent} per-horizon classifiers predict higher short- than long-term risk for $6.9\%$ of patients, the proposed model never ($0.0\%$; Cox, being proportional-hazards, is monotone too). Its behavior with event count also accounts for the real-data result below---retrained on growing subsamples and scored on a fixed test set, Cox leads at ${\approx}30$ events ($0.758$ vs.\ $0.618$, matching the $31$-event Mendeley cohort of Section~\ref{sec:results}-J) and the proposed model overtakes only above ${\approx}2000$ events (Fig.~\ref{fig:auroc}b): its extra parameters are a liability when events are few, an asset when many. We read this as \emph{consistent with} the real-data gap, not its proven cause (the endpoints also differ). Stratifying by predicted ten-year risk cleanly separates cumulative-incidence trajectories (Fig.~\ref{fig:cuminc}).

\begin{table}[!t]
\centering
\caption{AUROC by Horizon (Known-Status, \emph{Offer-Time} Feature Set). The 10-Year Differences Are Not Significant (See Paired CIs).}
\label{tab:auroc}
\begin{tabular}{@{}lccc@{}}
\toprule
Model & 1 month & 1 year & 10 year \\
\midrule
Logistic regression & 0.855 & 0.830 & 0.819 \\
Random forest & 0.851 & 0.825 & 0.818 \\
Gradient boosting (single) & 0.855 & 0.829 & 0.821 \\
Cox proportional hazards & 0.855 & 0.829 & 0.822 \\
Proposed (survival, MH) & 0.855 & 0.826 & 0.825 \\
\midrule
\multicolumn{4}{@{}p{7.6cm}}{Paired 10-y bootstrap: Proposed$-$GBM $+0.003$ $[-0.008,0.013]$; Proposed$-$Cox $+0.003$ $[-0.009,0.015]$} \\
\multicolumn{4}{@{}l}{Monotonicity violations: $0.0\%$ (proposed) vs.\ $6.9\%$ (independent)} \\
\bottomrule
\end{tabular}
\end{table}

\begin{figure}[!t]
\centering
\includegraphics[width=0.84\columnwidth]{out/fig_auroc.png}
\caption{Multi-horizon discrimination and the event-count crossover. (a) Known-status AUROC at each horizon for the fixed-horizon GBM, Cox PH, and the proposed multi-horizon survival model; the proposed model's ten-year edge over Cox falls within the paired bootstrap interval (inset). (b) Retrained on training sets of growing event count and scored on a fixed test set, Cox dominates when events are few---mirroring the $31$-event real cohort---and the proposed model overtakes only above ${\approx}2000$ events.}
\label{fig:auroc}
\end{figure}

\begin{figure}[!t]
\centering
\includegraphics[width=0.76\columnwidth]{out/fig7_cuminc.png}
\caption{Predicted cumulative-incidence trajectories by ten-year risk tertile; the model separates low-, medium-, and high-risk recipients.}
\label{fig:cuminc}
\end{figure}

\subsection{Calibration and Conformal Uncertainty}
For auditability we split the known-status one-year test set ($n=4519$) into disjoint isotonic-fit, conformal-calibration, and evaluation subsets, with the model fit on the separate training split. The raw proposed one-year risk is empirically well-calibrated (ECE $0.019$, Brier $0.132$); isotonic recalibration adds no held-out improvement (ECE $0.033$; Fig.~\ref{fig:calib})---an \emph{empirical} property, since the monotone-hazard construction constrains the ordering of risks across horizons, not their calibration. We apply split-conformal at level $\alpha=0.10$ to the better-calibrated raw probabilities (fixed before evaluation), measuring coverage on the disjoint, known-status one-year evaluation subset; size-two sets abstain for biopsy. Empirical marginal coverage is $91.0\%$ at the $90\%$ target with a $24.7\%$ abstention rate, and the retained singletons reach $0.852$ AUROC (Fig.~\ref{fig:conformal}). Because split-conformal is distribution-free, the score choice affects efficiency, not validity: the isotonic scores give similar coverage ($92.4\%$) but more abstention ($30.4\%$). Patients censored before the horizon lie outside the coverage claim; extending the guarantee to censored labels is future work.

\begin{figure}[!t]
\centering
\includegraphics[width=0.72\columnwidth]{out/fig_calibration.png}
\caption{One-year calibration of the proposed risk. The raw predictions already track the diagonal (ECE $0.019$ on held-out data); isotonic recalibration fit on a disjoint subset adds no held-out improvement.}
\label{fig:calib}
\end{figure}

\begin{figure}[!t]
\centering
\includegraphics[width=0.76\columnwidth]{out/fig10_conformal.png}
\caption{Conformal prediction. Empirical coverage tracks the target; the abstention rate is the cost paid for guaranteed coverage.}
\label{fig:conformal}
\end{figure}

\subsection{AMR Versus Tacrolimus Toxicity}
On the dysfunction cohort ($36\%$ toxicity), the gated classifier reaches $0.856$ AUROC and $0.770$ macro-F1. Uncertainty gating that refers the least-confident $21.5\%$ of cases for biopsy raises accuracy from $0.801$ to $0.861$ and macro-F1 from $0.770$ to $0.834$. AMR is recovered with $0.91$ recall; toxicity, the rarer and harder class, with $0.60$ (Fig.~\ref{fig:confusion})---precisely the cases the referral gate is designed to catch.

\begin{figure}[!t]
\centering
\includegraphics[width=0.6\columnwidth]{out/fig9_confusion.png}
\caption{Confusion matrix for the AMR-versus-toxicity differential on held-out dysfunction cases.}
\label{fig:confusion}
\end{figure}

\subsection{Decision-Focused Training and End-to-End Integration}
We evaluate the allocation module in isolation, and state its exact specification. This benchmark uses its \emph{own} analytic benefit target, distinct from the fairness benefit of Eq.~(4): it is built to carry nonlinear structure a linear learner cannot fit, so that the decision-focused comparison is meaningful. With per-pair features $\phi=[1,\,m,\,c,\,d_a,\,r_a,\,v]$ (standardized mismatch $m=|h_j-h_i|/6$, sensitization $c=(\mathrm{cPRA}_i-30)/30$, donor age $d_a$, recipient age $r_a$, vintage $v$), the target is
\begin{equation}
\begin{aligned}
B^{\mathrm{sp}}_{ji}=\mathrm{clip}_{[0.05,0.98]}\big(&\phi^\top\beta\;\underbrace{-0.35\,m\,\mathbb{1}[\mathrm{cPRA}_i{>}80]-0.25\,mc}_{\text{interactions}}\\[-2pt]
&\underbrace{-\,0.22\,m^2+0.15\,e^{-v_{\mathrm{yr}}/4}}_{\text{curvature, vintage}}\big),
\end{aligned}
\end{equation}
with $\beta=(0.80,-0.22,-0.10,-0.06,-0.04,0.05)$; clipping applies to the \emph{target} only. The two-stage and SPO+ predictors are both linear ridge regressions on the base map $\phi$, so both are misspecified against $B^{\mathrm{sp}}$; SPO+ minimizes the regret its predicted benefit induces, placing the optimizer inside training. This comparison runs on the \emph{utility-only} assignment---the fairness priority $s_i$ is held out of both arms---so it speaks to decision-focused allocation quality, not fair allocation.

Within the noise sweep, allocation regret rose monotonically with benefit-estimation error (Fig.~\ref{fig:regret}a; a sensitivity, not a law). We then trained two predictors of \emph{equal capacity} on identical donor batches, one with a mean-squared-error objective and one with the SPO+ decision-focused loss~\cite{spo}. Across ten seeds ($240$ held-out batches), SPO+ attains lower regret than two-stage ($1.84\%$ vs.\ $1.94\%$; paired reduction $0.094$ points, seed-clustered $95\%$ CI $[0.013,0.168]$; $9$ of $10$ seeds positive; Fig.~\ref{fig:regret}b)---and carries the signature of decision-focused learning: \emph{higher} prediction error ($0.0066$ vs.\ $0.0045$ MSE) yet lower regret, spending capacity on the pairs that decide the assignment. This advantage, however, is a \emph{misspecification} effect. Appending the target's four nonlinear terms to the feature map drops two-stage regret to $1.36\%$, and with the predictor now well specified the SPO+ gain falls to $0.019$ points, $95\%$ CI $[-0.045,0.086]$---indistinguishable from zero (Fig.~\ref{fig:regret}b). Getting the features right, not the training objective, is what closes the regret gap. This is the utility-only result; whether decision-focused training improves \emph{fair, feasible} allocation we now test directly.

\par\textbf{End-to-end integration.} We close the loop with the integrated objective of Eq.~(8): the fitted surrogate-benefit predictor feeds a fair, feasibility-constrained optimizer, and the two-stage and SPO+ predictors are trained and evaluated against the integrated regret with the NOTTO priority $s_i$ and the ABO/crossmatch feasible-edge set $\mathcal{F}$ (ABO-compatible, crossmatch-negative edges; the compatibility model is specified in Section~\ref{sec:results}-F) in \emph{both} arms (regret scored against the fair, feasible oracle). The utility-only advantage does not survive. At the baseline equity weight ($\lambda=0.12$, $40\times55$ pools, ten seeds), the paired reduction is $+0.037$ points with seed-clustered $95\%$ CI $[-0.030,0.102]$---indistinguishable from zero. Sweeping $\lambda$ from $0$ (utility-only) to $0.25$ leaves every interval straddling zero (Fig.~\ref{fig:regret}c); the gain reappears only at an extreme equity weight ($\lambda{=}0.5$: $+0.104$, CI $[0.029,0.183]$), where the priority term so dominates the objective that the misspecified predictor's residual errors resurface---not an operational setting. Across operating regimes (Fig.~\ref{fig:regret}d) the reduction is null except in the smallest, where with very scarce data SPO+ \emph{overfits} and loses significantly ($-0.265$ points, CI $[-0.623,-0.004]$); pushing the SPO+ training budget from $60$ to $300$ epochs does not open the gap, so the null is not an under-training artifact. A feasibility$\times$priority factorial localizes the effect (Table~\ref{tab:factorial}): under a common misspecified predictor the decision-focused advantage is significant only with \emph{both} constraints off ($+0.106$, CI $[0.029,0.181]$); each cell with a constraint active is individually indistinguishable from zero. We are deliberately careful here: the \emph{direct} between-cell contrasts are themselves inconclusive---priority-only minus unconstrained $-0.060$ (CI $[-0.179,0.051]$), feasibility-only $-0.053$ (CI $[-0.164,0.058]$), both $-0.069$ (CI $[-0.172,0.033]$), and the feasibility$\times$priority interaction $+0.045$ (CI $[-0.067,0.163]$)---so the factorial is \emph{consistent with} the gain being confined to the unconstrained regime but does not statistically establish that any single constraint reduces it. The operational reading is unchanged: at the baseline operating point we cannot distinguish the SPO+ and two-stage allocations, consistent with the feasible, high-priority matches being nearly forced (Section~\ref{sec:discussion}).

\begin{table}[!t]
\centering
\caption{Feasibility$\times$Priority Factorial Under a Common Misspecified Predictor: Paired Regret Reduction (Two-Stage $-$ SPO+) With Seed-Clustered 95\% CI Over Ten Seeds. The Decision-Focused Gain Is Significant Only When Both Operational Constraints Are Off.}
\label{tab:factorial}
\begin{tabular}{@{}llcc@{}}
\toprule
Feasibility & Priority & Reduction (pts) & 95\% CI \\
\midrule
off & off & $+0.106$ & $[+0.029,\,+0.181]$ \\
off & on  & $+0.046$ & $[-0.030,\,+0.114]$ \\
on  & off & $+0.053$ & $[-0.054,\,+0.156]$ \\
on  & on  & $+0.037$ & $[-0.030,\,+0.102]$ \\
\bottomrule
\end{tabular}
\end{table}

\begin{figure*}[!t]
\centering
\includegraphics[width=\textwidth]{out/fig_dfl_integration.png}
\caption{Decision-focused training, from allocation module to end-to-end integration. (a) Within the observed noise sweep, allocation regret rises monotonically with benefit-estimation error (a sensitivity, not a law). (b) In the \emph{utility-only} assignment at matched capacity, SPO+ achieves lower regret than two-stage for the base (misspecified) learner (seed-clustered $95\%$ CI excludes zero) but \emph{not} the enriched (well-specified) learner (CI crosses zero), localizing the gain to misspecification. (c) Under the \emph{integrated} fair, feasibility-constrained objective, the paired reduction (two-stage $-$ SPO+) has a $95\%$ CI that crosses zero at every equity weight up to $\lambda{=}0.25$ and reaches significance only at the extreme $\lambda{=}0.5$ ($*$). (d) Across operating regimes the reduction is null; only with very scarce data does SPO+ overfit and lose ($*$, CI excludes zero). Panels (c)--(d) use seed-clustered intervals over ten seeds.}
\label{fig:regret}
\end{figure*}

\par\textbf{Nonlinear capacity.} We also repeated the two extreme cells with a compact one-hidden-layer MLP ($16$ tanh units), trained both two-stage (MSE) and decision-focused (SPO+, with the optimizer-in-the-loop subgradient backpropagated through the network). Here the outcome differs from the linear case: at a fixed training budget (two-stage $700$ epochs, SPO+ warm-started for $150$) SPO+ lowers MLP regret in \emph{both} the unconstrained ($+0.361$, CI $[0.288,0.441]$) and the integrated ($+0.257$, CI $[0.198,0.320]$; ten of ten seeds) settings---a genuine advantage under this protocol. We suspect it is largely an optimization-budget effect: in the unconstrained setting the two-stage MLP's regret falls monotonically with budget ($0.86\%,\,0.43\%,\,0.25\%,\,0.21\%$ at $80$, $700$, $2000$, $4000$ epochs, three seeds) toward the SPO+ level ($\approx0.06\%$). But a declining curve is not convergence, we did not repeat the comparison in the integrated setting, and we report no continued-MSE control; we therefore state this as a \emph{hypothesis} and do \emph{not} claim the integration null extends to nonlinear learners. The decision-focused result is, in short, learner- and training-protocol-dependent, and settling whether the MLP gain survives a budget- and capacity-controlled comparison is future work.

\subsection{A Reduced-Form Survival-to-Allocation Study}
\label{sec:loop}
The main experiments use an analytic surrogate benefit. As a first step toward the survival-to-benefit coupling the architecture envisions (Fig.~\ref{fig:arch}), we feed a \emph{fitted survival-label proxy} into the fair, feasibility-constrained optimizer. The proxy is a ridge least-squares map fit to a binary ten-year graft-survival label ($\text{loss}\sim\text{Bernoulli}(1-B)$) with $35\%$ of labels removed completely at random. It is a simulated fixed-horizon survival signal, \emph{not} the paper's censoring-aware interval-hazard estimator: it models neither event times nor a survival likelihood, and its outputs are not probability-bounded. It therefore probes the coupling only in reduced form---connecting the full multi-horizon model with a proper likelihood and covariate-dependent censoring is future work. We compare a two-stage pipeline (fit the proxy, then plug it into the optimizer) against end-to-end decision-focused fine-tuning (SPO+). Because the end-to-end arm is handed the clean decision oracle while the two-stage arm sees only the noisy binary labels, this run carries \emph{privileged oracle information} for end-to-end and is not a neutral comparison. Under that caveat the end-to-end reduction is $+0.104$ points (seed-clustered $95\%$ CI $[0.027,0.181]$) unconstrained and $+0.010$ ($[-0.084,0.109]$) integrated; the integrated interval spans zero, so its sign is undetermined rather than a demonstrated equivalence. The reduced-form study is consistent with the unconstrained-only pattern seen elsewhere, but it does not establish that the main multi-horizon estimator has been connected to the allocator.

\subsection{Benchmarking the NOTTO Additive Priority}
We implement NOTTO's additive priority explicitly: a dialysis-time credit rewarding \emph{longer} waits plus group bonuses (pediatric largest), added to the benefit. This opposes the Eq.~(4) benefit, which rewards \emph{shorter} vintage---and that opposition is the price of fairness. We sweep the policy weight across three donor-scarcity regimes and eight independent pools each (Fig.~\ref{fig:fair}).

At the baseline donor-to-recipient ratio ($0.67$) the additive priority raises access for women ($n{\approx}121$; $0.68\to0.84$), highly sensitized candidates ($n{\approx}23$; $0.00\to0.38\pm0.11$), and pediatric candidates ($n{\approx}26$; $0.69\to0.99$), while utility falls only $0.85\pm0.16\%$; the mean dialysis vintage of allocated recipients rises from $3.3$ to $3.6$ years, confirming the credit reaches long-waiting patients. The highly sensitized group is the hardest: its low benefit resists reallocation, so its gain is partial and the most variable across pools---priority points alone cannot fully offset a biological disadvantage. The price of fairness tracks scarcity, widening to $1.24\pm0.20\%$ when organs are scarce (ratio $0.50$) and narrowing to $0.48\pm0.10\%$ when relatively abundant ($0.87$), because under scarcity every reassigned organ displaces a higher-benefit match (Fig.~\ref{fig:fair}a). Dispersion is the SD across the eight pools; the operating weight $\lambda^\star=0.24$ is a single illustrative point near the frontier knee, reused in the filtered scenario. ``Utility'' here is the Eq.~(4) surrogate benefit, and access figures are \emph{simulated} subgroup rates. As a sensitivity check we add ABO compatibility (O/B/A/AB $\approx 0.37/0.33/0.23/0.07$) and a crossmatch rule with incompatibility probability $0.02+0.9\,\mathrm{cPRA}_i/100$, so sensitized recipients lose most edges. This admits only $41\%$ of edges, yet nearly every donor still matches ($199.9$ of $200$), and the access gains are essentially unchanged (women $0.69\to0.84$, sensitized $0.02\to0.35$, pediatric $0.68\to0.99$; price of fairness $0.77\pm0.11\%$)---robust \emph{within this modeled feasibility scenario}, though only registry data could establish robustness to every compatibility structure.

\begin{figure*}[!t]
\centering
\includegraphics[width=0.64\textwidth]{out/fig_fairness.png}
\caption{NOTTO additive priority across donor-scarcity regimes (eight independent pools each). (a) Equity--utility frontier for scarce, baseline, and abundant donor-to-recipient ratios: raising the policy weight increases mean protected-group access while utility declines, and the frontier steepens as donors grow scarcer. (b) Per-group access at the baseline ratio under utility-only versus NOTTO priority, with group sizes ($n$) and across-pool error bars (price of fairness $0.85\%$).}
\label{fig:fair}
\end{figure*}

\subsection{Ablation and Feature Importance}
All rows use the offer-time \emph{GBM} as the reference ($0.829$ at one year); the proposed interval-hazard model---the primary predictor of Table~\ref{tab:auroc}---scores $0.826$ at one year, the small gap reflecting the two model classes rather than the feature set. The immunological block dominates one-year discrimination: from the offer-time baseline ($0.829$), dropping HLA, cPRA, and DSA collapses AUROC to $0.587$, while every other group moves it by at most $0.002$ (Table~\ref{tab:abl}). Adding \emph{measured tacrolimus trough} back---the post-transplant variable the offer-time model excludes (Section~II-H, Table~\ref{tab:timing})---does not improve one-year discrimination ($0.829$), and removing the germline CYP3A5 \emph{genotype} (retained at offer time) leaves it unchanged. No discrimination result depends on a post-transplant input.

\begin{table}[!t]
\centering
\caption{Ablation on the Offer-Time GBM: One-Year AUROC When a Feature Group Is Added or Removed.}
\label{tab:abl}
\begin{tabular}{@{}lcc@{}}
\toprule
Configuration & 1-year AUROC & $\Delta$ \\
\midrule
Offer-time model (primary) & 0.829 & -- \\
$+$ measured trough (research) & 0.829 & $+0.000$ \\
$-$ HLA, cPRA, DSA & 0.587 & $-0.242$ \\
$-$ CYP3A5 genotype & 0.829 & $+0.000$ \\
$-$ cold ischemia & 0.828 & $-0.001$ \\
$-$ dialysis vintage & 0.828 & $-0.001$ \\
\bottomrule
\end{tabular}
\end{table}

\subsection{Real-Data Validation on Open CKD Data}
To confirm that the detection pipeline runs end-to-end on genuinely open data, we applied it to the UCI Chronic Kidney Disease dataset~\cite{uci} (400 patients, Tamil Nadu, India; no licensing restriction, loaded through the \texttt{ucimlrepo} package). Five-fold cross-validated classifiers separate CKD from non-CKD almost perfectly (Table~\ref{tab:ckd}). We present this as a sanity check, not a hard benchmark: the cohort is small and strongly separable, so near-ceiling scores are expected. Reassuringly, the features the model relies on are the clinically established ones: hemoglobin, specific gravity, and albumin dominate the permutation importance (Fig.~\ref{fig:ckdimp}), with serum creatinine and packed cell volume close behind.

\begin{table}[!t]
\centering
\caption{CKD Detection on the Open UCI Dataset (5-Fold CV).}
\label{tab:ckd}
\begin{tabular}{@{}lccc@{}}
\toprule
Model & AUROC & Accuracy & F1 \\
\midrule
Logistic regression & 0.9999 & 0.995 & 0.996 \\
Random forest & 0.9999 & 0.995 & 0.996 \\
Gradient boosting & 0.9999 & 0.995 & 0.996 \\
\bottomrule
\end{tabular}
\end{table}

\begin{figure}[!t]
\centering
\includegraphics[width=0.8\columnwidth]{out/fig_ckd_importance.png}
\caption{Feature importance for CKD detection on the open UCI dataset. Hemoglobin, specific gravity, and albumin dominate, matching established clinical markers of chronic kidney disease.}
\label{fig:ckdimp}
\end{figure}

\subsection{Real-Data Graft-Rejection Prediction}
We next validated the rejection-prediction arm itself on real patients, using the openly licensed Mendeley kidney-allograft cohort~\cite{mendeley} (436 recipients from a UK transplant program; CC BY 4.0). Acute rejection (rate $8.9\%$) was predicted from pre-transplant features only---HLA mismatch, panel-reactive antibody, donor-specific antibody, crossmatch, ABO mismatch, induction regimen, and donor/recipient factors---with every post-transplant variable (graft function, infections, failure) withheld to prevent leakage. Baseline classifiers reach an AUROC of $0.79$ (random forest; Table~\ref{tab:transplant}), consistent with the pooled $\approx0.82$ of a recent meta-analysis~\cite{meta23}. High accuracy with low F1 reflects the class imbalance of a rare event, so AUROC is the appropriate summary. The pipeline thus transfers to real transplant data, while the allocation and fairness components stay synthetic because the cohort records completed transplants, not a live waiting list. The drivers are the expected ones---recipient age, donor-specific antibody, cold ischemia time, induction regimen, and HLA mismatch lead the permutation importance (Fig.~\ref{fig:tximp})---evidence the model learned transplant immunology, not dataset artifacts.

\begin{table}[!t]
\centering
\caption{Graft-Rejection Prediction on the Open Mendeley UK Cohort (436 Recipients, 5-Fold CV).}
\label{tab:transplant}
\begin{tabular}{@{}lccc@{}}
\toprule
Model & AUROC & Accuracy & F1 \\
\midrule
Logistic regression & 0.766 & 0.775 & 0.300 \\
Random forest & \textbf{0.788} & 0.901 & 0.295 \\
Gradient boosting & 0.707 & 0.908 & 0.259 \\
\bottomrule
\end{tabular}
\end{table}

\begin{figure}[!t]
\centering
\includegraphics[width=0.8\columnwidth]{out/fig_transplant_importance.png}
\caption{Feature importance for graft-rejection prediction on the real Mendeley cohort. Donor-specific antibody, cold ischemia time, induction regimen, and HLA mismatch rank among the top predictors, matching established immunological risk factors for rejection.}
\label{fig:tximp}
\end{figure}

\subsection{Real-Data Survival Validation and the Model-Capacity Trade-off}
We also tested the proposed survival estimator on a time-to-event endpoint using the same Mendeley cohort and the standard death-censored graft-loss definition ($31$ events among $434$ recipients; death with a functioning graft is treated as a censoring event, so competing-risk modeling of death is future work). Under five-fold evaluation---with imputation and standardization fit \emph{inside} each training fold, so no test information leaks---Harrell's C-index (Fig.~\ref{fig:survival}a) confirms real graft-loss signal: Cox leads at $0.735$ ($0.64$--$0.82$), the proposed model follows at $0.618$ ($0.51$--$0.72$), and the naive classifier trails at $0.599$ ($0.50$--$0.70$).

The proposed model does not beat Cox here, but the reason is instructive: with interval-specific hazards it estimates more parameters than Cox, so at $31$ events the parsimonious model wins---the unfavorable side of the capacity crossover of Section~\ref{sec:results}-A. A second public cohort confirms this: on $754$ Jakarta recipients~\cite{jakarta} with $86$ graft-loss events and the same features, the two are statistically level (proposed $0.557$ vs.\ Cox $0.540$; $+0.014$, $95\%$ CI $[-0.030,0.062]$; Fig.~\ref{fig:survival}a), as on the better-powered all-cause endpoint ($153$ events; $0.578$ vs.\ $0.566$). A tie at $86$ events with Cox ahead at $31$ is \emph{consistent with} the capacity account---the gap plausibly tracks event count rather than method, though the two cohorts differ in several respects, so we read it as a hypothesis, not an established cause. Absolute discrimination stays modest ($\sim$$0.56$) because this cohort's covariates are coarsely pre-dichotomized, while five-year calibration stays well-behaved (Fig.~\ref{fig:survival}b). We read the two cohorts as external confirmation that the features are predictive and the pipeline runs end-to-end on real data, while the clearest demonstration of the estimator's advantage remains the heavily censored regime it was designed for.

\begin{figure*}[!t]
\centering
\includegraphics[width=0.48\textwidth]{out/fig_survival.png}
\caption{Real-data survival validation across two independent cohorts (death-censored graft loss, five-fold out-of-fold). (a) Harrell C-index with bootstrap $95\%$ confidence intervals, proposed vs.\ Cox: on the Mendeley cohort ($31$ events) the parsimonious Cox model leads, but on the larger public Jakarta cohort ($86$ events) the two are statistically level---the gap tracks event count, not method. (b) Censoring-adjusted five-year calibration on the Jakarta cohort, by predicted-risk tertile.}
\label{fig:survival}
\end{figure*}

\section{Discussion}
\label{sec:discussion}
The novelty is not any single estimator---survival models, conformal prediction, and fair assignment all exist---but their composition into one pipeline, and the finding that governs it. In the unconstrained, utility-only assignment with a \emph{misspecified} predictor, decision-focused (SPO+) training helps: it carries higher prediction error than the two-stage predictor yet lower regret, by spending accuracy on the pairs that decide the assignment. The gain is fragile, though: it disappears once the feature map is corrected, and under the operational feasibility and equity constraints it is statistically indistinguishable from zero for the linear learner (Fig.~\ref{fig:regret}b--d, Table~\ref{tab:factorial}); the constraints already pin down the benefit-maximizing feasible match, so a simple two-stage predictor suffices. Two secondary analyses probe the edges rather than confirm the rule: a nonlinear (MLP) learner shows a decision-focused gain under our training budget that we cannot yet attribute cleanly to optimization, and a reduced-form survival-to-allocation proxy (Section~\ref{sec:loop}) reproduces only the unconstrained-only pattern. Neither overturns the linear result, and neither is claimed as confirmation. This is the paper's most useful and cautionary result---a decision-focused gain on a utility-only benchmark need not carry into a constrained deployment---so the guidance is to invest in the feasibility and equity layers and in calibration rather than in decision-focused training.

Clinically, the multi-horizon design matches how risk unfolds---early AMR and late chronic loss have different drivers a single horizon blurs---and calibrated abstention, referring ambiguous cases to biopsy, is a \emph{proposed}, simulation-only behavior a prospective study must test. For policy, soft NOTTO priorities turn an opaque scoring table into an auditable, tunable equity frontier.

\section{Limitations and Ethical Considerations}
The results are a proof of concept on synthetic data: they show the pipeline is well posed and internally consistent, not that it achieves these numbers clinically, so discrimination, calibration, and regret must be re-established on licensed registry data. The decision-focused comparison places the priority and feasible-edge set in both arms, but two couplings remain only partially addressed: Section~\ref{sec:loop} takes a first step on the survival-to-benefit link with a reduced-form survival-label proxy rather than the full interval-hazard model, and conformal abstention still acts only at the prediction stage---connecting the multi-horizon estimator with a proper likelihood and folding abstention into a single trained pipeline are future work. The integration null is established for the linear predictor; for a one-hidden-layer MLP the decision-focused gain persists under our training budget, so whether the result extends to higher-capacity learners is open. Fairness is a soft priority for three groups, biasing access rather than guaranteeing a floor; hard reserved-share constraints would trade off differently. The AMR-versus-toxicity differential, the conformal referral, and every allocation and fairness result are evaluated \emph{only in simulation}.

Organ allocation is high-stakes and contested; any deployed model must be externally and prospectively validated, auditable, and clinician-governed rather than autonomous. Data governance must comply with India's DPDP Act, 2023, motivating the federated design that keeps patient data local. Protected-attribute use is itself a policy choice; our role is to make the trade-off explicit and tunable, not to set it.

\section{Conclusion and Future Work}
We presented a modular framework that couples multi-horizon graft-rejection prediction, calibrated uncertainty, and a NOTTO-inspired fairness-constrained allocator, tested on synthetic data and, for the prediction arms, on three open real datasets. Its most transferable result is honest and negative: decision-focused training helps only a misspecified predictor in the unconstrained allocation, and under operational feasibility and equity constraints it is statistically indistinguishable from a simple two-stage predictor for the linear learner we study; whether this extends to higher-capacity learners our nonlinear check does not settle---so those constraints and good calibration, not the decision-focused objective, govern allocation quality. The next steps are to substitute fitted survival probabilities and conformal abstention for the benefit surrogate in one trained pipeline, and to re-estimate on registry (SRTR/OPTN) data under a prospective, clinician-governed study. Until then this is a proof of concept---Cox leads at $31$ real events but draws level at $86$, and the allocation results measure \emph{surrogate} utility---and none of it is clinical guidance.

\begin{thebibliography}{99}
\setlength{\itemsep}{0pt plus 0.2pt}
\bibitem{notto} M.\ Sai Gopal, ``One nation, one policy: major overhaul for organ donation,'' \textit{Telangana Today}, 18 Dec.\ 2025. [Online]. Available: \url{https://telanganatoday.com/one-nation-one-policy-major-overhaul-for-organ-donation}
\bibitem{nottoadv} Press Trust of India, ``Give priority to women, close kin of deceased donors for organ transplant: Govt to states,'' \textit{ThePrint}, 2025. [Online]. Available: \url{https://theprint.in/india/give-priority-to-women-close-kin-of-deceased-donors-for-organ-transplant-govt-to-states/2718068}
\bibitem{reva} REVA RACE Lab, ``Deep-learning-based deceased-donor kidney allocation model for India (NOTTO),'' [Online]. Available: \url{https://race.reva.edu.in/race-lab/deep-learning-based-deceased-donor-kidney-allocation}
\bibitem{oneindia} S.\ Banerjee \textit{et al.}, ``Is India ready for `One India One Allocation' in deceased-donor kidney transplant?'' \textit{Indian J. Transplantation}, 2025. [Online]. Available: \url{https://www.citedrive.com/en/discovery/is-india-ready-for-one-india-one-allocation-policy-in-deceased-donor-kidney-transplant/}
\bibitem{luck} M.\ Luck, T.\ Sylvain, H.\ Cardinal, A.\ Lodi, and Y.\ Bengio, ``Deep learning for patient-specific kidney graft survival analysis,'' \textit{arXiv:1705.10245}, 2017. [Online]. Available: \url{https://arxiv.org/abs/1705.10245}
\bibitem{bae} S.\ Bae \textit{et al.}, ``Machine learning to predict transplant outcomes: helpful or hype?'' \textit{Transplant International}, vol.\ 33, no.\ 11, pp.\ 1472--1480, 2020.
\bibitem{meta23} B.\ Ravindhran, P.\ Chandak, N.\ Schafer, K.\ Kundalia, W.\ Hwang, S.\ Antoniadis, U.\ M.\ Haroon, and R.\ H.\ Zakri, ``Machine learning models in predicting graft survival in kidney transplantation: meta-analysis,'' \textit{BJS Open}, vol.\ 7, no.\ 2, zrad011, 2023, doi:10.1093/bjsopen/zrad011.
\bibitem{dibello} F.\ Di Bello \textit{et al.}, ``The role of artificial intelligence in predicting graft survival in kidney transplantation: a systematic review,'' \textit{Minerva Urol.\ Nephrol.}, vol.\ 78, no.\ 1, pp.\ 53--62, 2026.
\bibitem{klundert} J.\ van de Klundert, F.\ P\'erez-Galarce, M.\ Olivares, L.\ Pengel, and A.\ de Weerd, ``The comparative performance of models predicting patient and graft survival after kidney transplantation: a systematic review,'' 2024.
\bibitem{tfn} A.\ Kumar, S.\ Rauch, M.\ Cypko, M.\ Naik, and M.-P.\ Schapranow, ``Temporal Fusion Nexus: a task-agnostic multimodal embedding model for post-kidney-transplant care,'' \textit{arXiv:2601.08503}, 2026.
\bibitem{belcic} T.\ Bel\v{c}i\v{c} Miki\v{c} and M.\ Arnol, ``The use of machine learning in the diagnosis of kidney allograft rejection,'' \textit{Diagnostics}, vol.\ 14, no.\ 22, 2482, 2024.
\bibitem{almoussawy} M.\ Al Moussawy, Z.\ S.\ Lakkis, Z.\ A.\ Ansari, A.\ R.\ Cherukuri, and K.\ I.\ Abou-Daya, ``The transformative potential of artificial intelligence in solid organ transplantation,'' \textit{Front.\ Transplant.}, vol.\ 3, 1361491, 2024.
\bibitem{wsi} Y.\ Ye, L.\ Xia, S.\ Yang, Y.\ Luo, Z.\ Tang, Y.\ Li, L.\ Han, H.\ Xie, Y.\ Ren, and N.\ Na, ``Deep learning-enabled classification of kidney allograft rejection on whole slide histopathologic images,'' \textit{Front.\ Immunol.}, vol.\ 15, 1438247, 2024, doi:10.3389/fimmu.2024.1438247.
\bibitem{divard} G.\ Divard, M.\ Raynaud, V.\ S.\ Tatapudi \textit{et al.}, ``Comparison of artificial intelligence and human-based prediction and stratification of the risk of long-term kidney allograft failure,'' \textit{Commun.\ Med.}, vol.\ 2, no.\ 1, 2022, doi:10.1038/s43856-022-00201-9.
\bibitem{thong} C.\ Thongprayoon \textit{et al.}, ``Feature importance of acute rejection among Black kidney transplant recipients,'' \textit{Medicines}, vol.\ 8, no.\ 11, 66, 2021.
\bibitem{offei} E.\ Offei \textit{et al.}, ``Does eplet load and electrostatic mismatch score matter? A machine-learning approach,'' \textit{Transplantology}, vol.\ 6, no.\ 1, 6, 2025.
\bibitem{tacro} J.\ Min, Q.\ Li, W.\ Lai, Y.\ Lu, X.\ Wang, and G.\ Chen, ``Clinical-oriented tacrolimus dosing via genetic algorithm and deep forest,'' \textit{Front.\ Pharmacol.}, vol.\ 16, 1656197, 2025.
\bibitem{zhang} M.\ Zhang, X.\ Dai, and L.\ Li, ``Fairness-aware organ exchange and kidney paired donation,'' \textit{arXiv:2503.06431}, 2025.
\bibitem{carvalho} M.\ Carvalho \textit{et al.}, ``Penalties and rewards for fair learning in paired kidney exchange programs,'' \textit{arXiv:2309.13421}, 2023.
\bibitem{spo} A.\ N.\ Elmachtoub and P.\ Grigas, ``Smart `predict, then optimize','' \textit{Management Science}, vol.\ 68, no.\ 1, pp.\ 9--26, 2022.
\bibitem{bert} D.\ Bertsimas, V.\ F.\ Farias, and N.\ Trichakis, ``Fairness, efficiency, and flexibility in organ allocation for kidney transplantation,'' \textit{Oper.\ Res.}, vol.\ 61, no.\ 1, pp.\ 73--87, 2013, doi:10.1287/opre.1120.1138.
\bibitem{uci} L.\ Rubini, P.\ Soundarapandian, and P.\ Eswaran, ``Chronic Kidney Disease,'' UCI Machine Learning Repository, 2015, doi:10.24432/C5G020. [Online]. Available: \url{https://archive.ics.uci.edu/dataset/336/chronic+kidney+disease}
\bibitem{mendeley} C.\ A.\ Chukwu, ``Outcomes of basiliximab and alemtuzumab induction in kidney allograft recipients with matched immunological risk profile: a retrospective matched cohort study,'' \textit{Mendeley Data}, V1, 2025, doi:10.17632/3vdszzmz2n.1. [Online]. Available: \url{https://data.mendeley.com/datasets/3vdszzmz2n/1}
\bibitem{jakarta} ``Long-term outcomes and prognostic factors in kidney transplant recipients in Jakarta, Indonesia: a cohort study,'' \textit{Dryad Dataset}, 2022, doi:10.5061/dryad.41ns1rnfr. [Online]. Available: \url{https://datadryad.org/dataset/doi:10.5061/dryad.41ns1rnfr}
\end{thebibliography}

\end{document}
